%%%%%%%%%%%%%%%%%%%%%%%%%%%%%%%%%%%
%%%This is the LaTeX ARTICLE template for ITU journals
%%%%%%%%%%%%%%%%%%%%%%%%%%%%%%%%%%%
\documentclass[twoside,twocolumn,9pt]{article}
\usepackage{itujournalstyle}

%%%Settings specific to the ITU journal


%%%Enable the following three lines for ITU JWCC 
\definecolor{jour_color}{rgb}{0.09, 0.24, 0.49} 
\newcommand{\Jname}{JWCC}
\newcommand{\refstyle}{ieee}



%%%Please enter below first page info, such as title, abstract, 
%%%author names, affiliations, etc.

%%Article title goes here instead of the text "Single Space Title"
\newcommand{\TitleJournal}{Single Space Title of your Paper in Arial 15 pt Goes Here}

%%%Türkçe başlık burada yer almalıdır
\newcommand{\TurkceBaslik}{Başlıkta İkinci Satıra Geçen Satırlar Arasında Boşluk Olmamalıdır}

%%%The abstract goes here instead of the text "A single paragraph of ..."
\newcommand{\AbstractJournal}{\justify{A single paragraph of about 250-300 words maximum. For research articles, abstracts should give a pertinent overview of the work. Lorem ipsum dolor sit amet, consectetur adipiscing elit, sed do eiusmod tempor incididunt ut labore et dolore magna aliqua. Rhoncus dolor purus non enim praesent elementum facilisis leo vel. Neque sodales ut etiam sit amet nisl. Non quam lacus suspendisse faucibus interdum posuere. Pellentesque elit ullamcorper dignissim cras. Ornare quam viverra orci sagittis. Sapien efwwfewfwefget mi proin sed libero enim. Mi quis hendrerit dolor magna eget. Euismod quis viverra nibh cras pulvinar mattis nunc. Fringilla ut morbi tincidunt augue interdum velit euismod in pellentesque. Sed arcu non odio euismod lacinia at.}} 

%%%Türkçe Özet aşağıda verilmelidir
\newcommand{\Ozet}{\justify{Maksimum 250-300 kelime arasında olmalı ve buraya yazı fontu (Arial 9) yazılmalıdır. Metin yazılırken bu taslak indirilip üzerinden devam edilmelidfwefefwqefir. Lorem ipsum dolor sit amet, consectetur adipiscing elit, sed do eiusmod tempor incididunt ut labore et dolore magna aliqua. Rhoncus dolor purus non enim praesent elementum facilisis leo vel. Neque sodales ut etiam sit amet nisl. Non quam lacus suspendisse faucibus interdum posuere. Pellentesque elit ullamcorper dignissim cras. Ornare quam viverra orci sagittis. Sapien eget mi proin sed libero enim. Mi quis hendrerit dolor magna eget.}} 
\newcommand{\Keywords}{First, second, third, max 5 keywords}
\newcommand{\AnahtarKelimeler}{Birinci, ikinci, üçüncü, maks 5 kelime}


%Authors' names and affiliations go below instead of "Author1", %"Author2", ... and "Affiliation1", "Affiliation2" etc.
\newcommand{\AuthorA}{Author1}
\newcommand{\AuthorB}{Author2}
\newcommand{\AuthorC}{Author3}
%\newcommand{\AuthorC}{Author4}
%...

\newcommand{\AdresA}{Affiliation1}
\newcommand{\AdresB}{Affiliation2}
\newcommand{\AdresC}{Affiliation3}
%\newcommand{\AdresC}{Affiliation4}
%...

%Enter orcid IDs for the authors
\newcommand{\orcidauthorA}{0000-0000-0000-000X} 
\newcommand{\orcidauthorB}{0000-0000-0000-000X} 
\newcommand{\orcidauthorC}{0000-0000-0000-000X} 

%%%Reference info, keep as is
 \usepackage[
    backend=biber,
    style=\refstyle,
  ]{biblatex}
  \usepackage{csquotes}
\addbibresource{sample.bib}

\begin{document}
\justifying{
\input{headfoot.tex}
%%%Main Text%%%%
\section{Introduction}
The introduction should briefly place the study in a broad context and highlight why it is important. It should define the purpose of the work and its significance. The current state of the research field should be carefully reviewed and key publications cited (\cite{einstein}).

Finally, briefly mention the main aim of the work and highlight the principal conclusions.
\section{Material}
This section may be divided by subheadings. 
\subsection{Materials}
\subsection{Processing techniques}

Main headings should be informative, use 9 point Arial font and capitalize the first letter of the all words.

For Level 2 headings, use justified 9-point Arial font and capitalize the first letter

Section headings can be typeset with and without numbers(\cite{knuthwebsite}).

\subsubsection{This is the subsubsection style.~~} These headings should end in a full point.  

\section{Graphics and tables}
\subsection{Graphics}
Graphics should be inserted on the page where they are first mentioned (unless they are equations, which appear in the flow of the text) (\cite{dirac}).

\begin{figure}[ht]
\centering
  \includegraphics[height=4cm]{fig_example.jpg}
  \caption{An example figure caption.}
  \label{fgr:example}
\end{figure}

\subsection{Tables}
Tables typeset in RSC house style do not include vertical lines. Table footnote symbols are lower-case italic letters and are typeset at the bottom of the table. Table captions do not end in a full point (\cite{dirac,einstein,knuth-fa}).

\begin{table}[ht]
\small
  \caption{\ An example of a caption to accompany a table \textendash\ table captions do not end in a full point}
  \label{tbl:example1}
  \begin{tabular*}{0.48\textwidth}{@{\extracolsep{\fill}}lll}
    \hline
    Header one (units) & Header two & Header three \\
    \hline
    1 & 2 & 3 \\
    4 & 5 & 6 \\
    7 & 8 & 9 \\
    10 & 11 & 12 \\
    \hline
  \end{tabular*}
\end{table}

Adding notes to tables can be complicated.  Perhaps the easiest method is to generate these manually.

\section{Equations}

Equations can be typeset inline \textit{e.g.}\ $ y = mx + c$ or displayed with and without numbers:
%The \balance command can be used to balance the columns on the final page if desired. It should be placed anywhere within the first column of the last page.
\balance
 \[ A = \pi r^2 \]

\begin{equation}
  \frac{\gamma}{\epsilon x} r^2 = 2r
\end{equation}

You can also put lists into the text. You can have bulleted or numbered lists of almost any kind. 

Lists are easy to create:
\begin{itemize}
%for numbered list use \begin{enumerate} 
  \item List entries start with the \verb|\item| command.
  \item Individual entries are indicated with a black dot, a so-called bullet.
  \item The text in the entries may be of any length.
\end{itemize}

The \texttt{mhchem} package can also be used so that formulae are easy to input: \texttt{\textbackslash ce\{H2SO4\}} gives \ce{H2SO4}. 

\section{Conclusion}
The conclusions section should come in this section at the end of the article, before the Author Contributions or References.

\section*{Author Contributions}
We encourage authors to include author contributions. 

\section*{Acknowledgements}
The Acknowledgements come at the end of an article after Conflicts of interest and before the references.

%%%END OF MAIN TEXT%%%



%you can change the title of 'References' to for instance %Reference List using the following command:
%\renewcommand\refname{List of References}

%%%References%%%
\printbibliography %Prints bibliography
}
\end{document}
